\documentclass[10pt,twocolumn,a4paper]{article}

% ==============================================================================
% PAQUETES Y CONFIGURACIÓN EDITORIAL (DISEÑO CON ENCABEZADOS FANCYHDR)
% ==============================================================================
\usepackage[utf8]{inputenc}
\usepackage[T1]{fontenc}
\usepackage[spanish,es-tabla]{babel}
\usepackage{amsmath,amssymb,amsfonts}
\usepackage{graphicx}
\usepackage{booktabs}
\usepackage{geometry}
\usepackage{hyperref}
\usepackage{cite}
\usepackage{caption}
\usepackage{float}
\usepackage{xcolor}
\usepackage{fancyhdr}
\usepackage{microtype}

\geometry{top=2.0cm, bottom=2.0cm, left=1.55cm, right=1.55cm, columnsep=0.60cm, headheight=14pt}

\definecolor{tharosblue}{RGB}{18, 65, 145}
\hypersetup{
    colorlinks=true,
    linkcolor=tharosblue,
    citecolor=tharosblue,
    urlcolor=tharosblue,
    pdftitle={Priones de Carbono: Nucleación Conformacional Dirigida y Dimerización Cuántica de Diamantoides Dopados a Baja Temperatura},
    pdfauthor={César Leonardo Rodríguez García (Dark Horse / NocWoeder165)}
}

% Cabezales de página elegantes (estilo artículo científico)
\pagestyle{fancy}
\fancyhf{}
\lhead{\footnotesize\textit{Pre-print -- Proyecto Tharos (DOI: 10.5281/zenodo.23427695)}}
\rhead{\footnotesize\textit{C.~L.~Rodríguez García -- Dark Horse / NocWoeder165 (2026)}}
\cfoot{\thepage}
\renewcommand{\headrulewidth}{0.4pt}

\captionsetup{font=footnotesize, labelfont=bf}

\begin{document}

% ==============================================================================
% ENCABEZADO Y RESUMEN A ANCHO COMPLETO (ESTILO PRE-PRINT CLÁSICO)
% ==============================================================================
\twocolumn[
  \begin{@twocolumnfalse}
    \vspace{-0.5cm}
    {\centering
      {\LARGE \textbf{Priones de Carbono: Nucleación Conformacional Dirigida y Dimerización Cuántica de Diamantoides Dopados a Baja Temperatura} \par}
      \vspace{0.85em}
      {\large \textbf{César Leonardo Rodríguez García} \par}
      \vspace{0.3em}
      {\small \textit{Universidad Tecnológica de Nayarit (UTNAY), Xalisco, Nayarit, México} \par}
      \vspace{0.15em}
      {\footnotesize \textit{Proyecto Tharos} $\cdot$ Contacto: \texttt{sem-340007@utnay.edu.mx} \par}
      \vspace{0.25em}
      {\footnotesize 18 Octubre 2026 $\cdot$ \textbf{Pre-print DOI:} \href{https://doi.org/10.5281/zenodo.23427695}{\texttt{10.5281/zenodo.23427695}} \par}
      \vspace{0.65em}
    }
    \hrule
    \vspace{0.55em}
    \noindent\textbf{\large Resumen}
    \vspace{0.35em}

    \small
    Durante siete décadas se ha asumido que la síntesis de diamante exige regímenes de fuerza bruta: presiones de gigapascales (HPHT) o plasmas térmicos a $>900\text{ }^\circ\text{C}$ (MPCVD). Dichas condiciones destruyen sustratos microelectrónicos e inducen segregación de dopantes. Sin embargo, al observar la molécula de adamantano ($\text{C}_{10}\text{H}_{16}$), emerge una contradicción fundamental: la celda unitaria $sp^3$ del diamante ya existe a temperatura ambiente sin necesidad de compresión externa. En este estudio computacional se formula y audita en dos niveles de teoría la hipótesis de los \textbf{``Priones de Carbono''}: el uso de jaulas diamantoides pre-estructuradas para transferir orden tetraédrico en alto vacío ($\sim 10^{-6}\text{ Torr}$) a baja temperatura ($< 250\text{ }^\circ\text{C}$). En una primera fase de exploración con el potencial interatómico fundacional MACE-MP-0 acoplado a dispersión D3 (\texttt{float64}), el escaneo relajado entre radicales 1-adamantilo ($\text{C}_{10}\text{H}_{15}^\bullet$) y derivados trivalentes neutros sustituidos con boro ($\text{C}_9\text{BH}_{14}^\bullet$) y fósforo ($\text{C}_9\text{PH}_{14}^\bullet$) exhibió crestas aparentes de $0.251\text{ a }0.642\text{ eV}$ entre $2.60\text{ y }2.44\text{ \AA}$. Para resolver la naturaleza electrónica de dicha región, se ejecutó un protocolo de 27 cálculos \textit{ab initio} de espín abierto con simetría rota (\texttt{BS-UKS} $\omega\text{B97X-D3/def2-SVP}$, proyección de Yamaguchi \texttt{AP-UKS} y estado triplete \texttt{HS-UKS}) sobre geometrías relajadas en \texttt{UKS B97-3c} en ORCA 6.1.1. La auditoría cuántica revela que la cresta de MACE-MP-0 (modelo ciego al espín) seguía la rama repulsiva del estado Triplete ($S=1, \uparrow\uparrow$, $\Delta E_{\text{HS}} = +0.18\text{ a }+0.58\text{ eV}$ por exclusión de Pauli). En contraste, el canal Singlete Antiferromagnético ($M_S=0, \uparrow\downarrow$) desciende sin barrera energética conforme el valor esperado $\langle \hat{S}^2 \rangle_{\text{BS}}$ colapsa de $\sim 0.95$ a $0.000$, estabilizando pozos covalentes profundos ($\Delta E = -3.921\text{ eV}$ en $\text{C}_{20}\text{H}_{30}$, $-3.651\text{ eV}$ en $\text{C}_{19}\text{H}_{29}\text{B}$ y $-3.727\text{ eV}$ en $\text{C}_{19}\text{H}_{29}\text{P}$). La optimización \texttt{UKS B97-3c} reproduce la distancia experimental por rayos X del 1,1'-biadamantilo ($1.578\text{ \AA}$) con una precisión de $1.5771\text{ \AA}$ ($-0.06\%$). Estadísticamente, en un haz térmico no polarizado, el $25\%$ de las colisiones bimoleculares ocurre por el canal singlete sin barrera, mientras que el $75\%$ triplete experimenta rebote de Pauli. Estos resultados acotan con rigor \textit{ab initio} el paso elemental de dimerización aislada, sin demostrar aún el crecimiento cristalino continuo ni la activación eléctrica del dopaje.

    \vspace{0.45em}
    \noindent\textbf{Palabras clave:} Carbon prions; Diamondoids; Broken-symmetry UKS DFT; ORCA 6.1.1; $\omega\text{B97X-D3}$; MACE-MP-0; Spin statistics.
    \vspace{0.55em}
    \hrule
    \vspace{1.0em}
  \end{@twocolumnfalse}
]

% ==============================================================================
% 1. EL DOGMA DE LA FUERZA BRUTA
% ==============================================================================
\section{La Paradoja de la Fuerza Bruta}
El diamante monocristalino posee las propiedades físicas más extremas de la materia condensada: brecha prohibida de $5.47\text{ eV}$, campo de ruptura superior a $10\text{ MV/cm}$ y conductividad térmica por encima de $2000\text{ W/m}\cdot\text{K}$ \cite{Koizumi2001}. Pese a ello, la microelectrónica moderna continúa atada al silicio. ¿Por qué no existen procesadores cotidianos sobre obleas de diamante semiconductor?

La respuesta reside en el método de fabricación. Tradicionalmente se ha creído que para ordenar carbono en una red tetraédrica ($sp^3$) solo existen dos caminos: aplastar fases grafíticas bajo presiones de $5\text{ a }17\text{ GPa}$ (HPHT / DAC) \cite{Suzuki2024}, o encender plasmas de microondas a $700\text{--}900\text{ }^\circ\text{C}$ (MPCVD, desarrollado originalmente por Kamo et al.\ en 1983 \cite{Kamo1983}) para fragmentar metano ($\text{CH}_4$) en un bombardeo caótico de radicales atómicos. Se gasta una cantidad desproporcionada de energía para destruir enlaces previos y luego intentar que los átomos caigan en los sitios correctos. Cuando se intenta introducir boro (aceptor tipo $p$, energía de ionización $\sim 0.37\text{ eV}$) o fósforo (donador profundo tipo $n$, $\sim 0.57\text{ eV}$) bajo ese régimen térmico violento, la red expulsa las impurezas hacia la superficie o colapsa hacia defectos grafíticos ($sp^2$) \cite{Koizumi2001}.

Sin embargo, basta observar una molécula de adamantano ($\text{C}_{10}\text{H}_{16}$) para notar que el problema de la hibridación local ya está resuelto a temperatura ambiente \cite{Schleyer1957,Dahl2003}. El adamantano no se parece al diamante: \textit{es} un fragmento estequiométrico de diez carbonos de la red cristalina del diamante estabilizado por dieciséis hidrógenos periféricos, capaz de autoensamblarse en monocapas ordenadas en vacío \cite{Yang2007}. En 2025, Fu, Nakamuro y Nakamura confirmaron en \textit{Science} que al arrancar hidrógenos de cristales de adamantano con electrones de $80\text{--}200\text{ keV}$, las jaulas se sueldan espontáneamente formando nanodiamantes esféricos a temperaturas de $100\text{--}296\text{ K}$ \cite{Fu2025}.

Si la geometría $sp^3$ ya viene empaquetada de fábrica dentro del monómero, la pregunta cambia por completo: ¿realmente se necesita presión hidrostática o plasma térmico para unir dos jaulas, o basta con activar un vértice y dejar que la propia complementariedad geométrica y de espín haga el trabajo?

% ==============================================================================
% 2. EL PRINCIPIO PRIÓNICO Y PROTOCOLO EN DOS NIVELES (MACE + ORCA 6.1.1)
% ==============================================================================
\section{El Principio Priónico y Metodología de Doble Nivel}
En biofísica, una proteína priónica actúa como un molde conformacional que induce el acoplamiento hacia un estado termodinámicamente profundo mediante reconocimiento de interfaz. Trasladado como analogía heurística a la química del estado sólido, el concepto de \textbf{``Priones de Carbono''} plantea que una jaula diamantoide activada expone un molde tetraédrico rígido donde el acoplamiento frontal sustituye la compresión mecánica por información estérica pre-codificada, en línea con la reducción de barreras de nucleación medida en semillas diamantoides \cite{Gebbie2018}.

Para evaluar con rigor el paso elemental de esta hipótesis (la dimerización bimolecular aislada en vacío) y evitar artefactos de potenciales empíricos o redes neuronales ciegas al espín, se implementó una arquitectura de validación en dos niveles:

\subsection{Nivel 1: Exploración Preliminar con MACE-MP-0 + D3}
Se modelaron inicialmente tres reacciones frontales en el entorno ASE (\texttt{PyTorch 2.11.0}, \texttt{ASE 3.29.0}, \texttt{MACE 0.3.16}) empleando el potencial neuronal equivariante \textbf{MACE-MP-0} (\textit{small}) \cite{Batatia2024} en precisión doble (\texttt{float64}) acoplado a dispersión de Grimme D3 (\texttt{bj}) obligatoria \cite{Grimme2010}: (i) acoplamiento intrínseco ($\text{C}_{10}\text{H}_{15}^\bullet + \text{C}_{10}\text{H}_{15}^\bullet \to \text{C}_{20}\text{H}_{30}$); (ii) acoplamiento con un radical 1-bora-adamantilo trivalente neutro ($\text{C}_9\text{BH}_{14}^\bullet \to \text{C}_{19}\text{H}_{29}\text{B}$); y (iii) acoplamiento con 1-fosfa-adamantilo trivalente neutro ($\text{C}_9\text{PH}_{14}^\bullet \to \text{C}_{19}\text{H}_{29}\text{P}$). Partiendo del adamantano cristalográfico (PubChem CID 9238), las jaulas se alinearon en conformación alternada (\textit{staggered}, $60^\circ$) mediante rotaciones exactas en $\text{SO}(3)$, realizando un barrido relajado en 13 pasos ($3.40 \to 1.48\text{ \AA}$, $F_{\text{max}} < 0.05\text{ eV/\AA}$) y una relajación final libre en el pozo ($F_{\text{max}} < 0.02\text{ eV/\AA}$).

\subsection{Nivel 2: Auditoría Ab Initio de Espín Abierto (ORCA 6.1.1)}
Dado que MACE-MP-0 fue entrenado sobre cálculos PBE de capa cerrada y carece de grados explícitos de espín para describir el acoplamiento birradicalario ($S_1 = 1/2 \otimes S_2 = 1/2 = S_0 \oplus T_1$), se ejecutó un lote de 27 cálculos \textit{ab initio} de capa abierta en \textbf{ORCA 6.1.1} \cite{Neese2022} empleando el funcional híbrido de largo alcance corregido por dispersión \textbf{UKS $\omega\text{B97X-D3/def2-SVP}$} (\texttt{def2/J RIJCOSX TightSCF}) \cite{Chai2013} y geometrías optimizadas con el método compuesto \textbf{UKS B97-3c} \cite{Brandenburg2018}. Para garantizar que ningún punto mezcle geometrías de distintos métodos, se definieron dos esquemas estrictamente homogeneizados respecto al cero termodinámico de radicales dobletes aislados ($S=1/2$, $\langle \hat{S}^2 \rangle \approx 0.755\text{--}0.772$):
\begin{equation}
    E_\infty = E_{\text{radical\_1}}^\circ + E_{\text{radical\_2}}^\circ
\end{equation}
\begin{itemize}
    \item \textbf{Opción A (\texttt{BS-UKS} $\omega\text{B97X-D3/def2-SVP}$ // \texttt{MACE-MP-0+D3}):} Evalúa las geometrías relajadas del barrido MACE contra los monómeros en geometría MACE ($E_\infty^{\text{A}} = -779.379994\text{ }E_{\text{h}}$ para $\text{C}_{20}\text{H}_{30}$; $-765.529623\text{ }E_{\text{h}}$ para $\text{C}_{19}\text{H}_{29}\text{B}$; y $-1081.940273\text{ }E_{\text{h}}$ para $\text{C}_{19}\text{H}_{29}\text{P}$).
    \item \textbf{Opción B (\texttt{BS-UKS} $\omega\text{B97X-D3/def2-SVP}$ // $C_3$-\texttt{UKS-B97-3c}):} Evalúa la aproximación coaxial a lo largo del eje $C_3$ de los monómeros optimizados en \texttt{UKS B97-3c} (y el dímero libremente optimizado en \texttt{UKS B97-3c}) contra los monómeros \texttt{UKS B97-3c} ($E_\infty^{\text{B}} = -779.401978\text{ }E_{\text{h}}$, $-765.551630\text{ }E_{\text{h}}$ y $-1081.973354\text{ }E_{\text{h}}$; con estabilizaciones de relajación monomérica $\Delta E_{\text{relaj}}^{\text{MACE}\to\text{B97-3c}}$ de $-0.598$, $-0.599$ y $-0.900\text{ eV}$, respectivamente).
\end{itemize}
En cada distancia inter-jaula ($3.40, 2.92, 2.60, 2.44, 2.12\text{ \AA}$ y $d_{\text{eq}}$) se resolvió la directiva \texttt{BrokenSym 1,1}, obteniendo simultáneamente el estado Triplete de alto espín ($E_{\text{HS}}, S=1, \uparrow\uparrow, \langle \hat{S}^2 \rangle_{\text{HS}} \approx 2.01$), el Singlete birradicalario de simetría rota ($E_{\text{BS}}, M_S=0, \uparrow\downarrow$) y la energía del estado Singlete puro corregida por contaminación de espín mediante la proyección aproximada de Yamaguchi ($\text{AP-UKS}$) \cite{Yamaguchi1988}:
\begin{equation}
    E_{\text{AP}} = E_{\text{BS}} + \frac{\langle \hat{S}^2 \rangle_{\text{BS}}}{\langle \hat{S}^2 \rangle_{\text{HS}} - \langle \hat{S}^2 \rangle_{\text{BS}}} \left(E_{\text{BS}} - E_{\text{HS}}\right)
\end{equation}

% ==============================================================================
% FIGURA 1 A ANCHO COMPLETO (EXPLORACIÓN MACE-MP-0)
% ==============================================================================
\begin{figure*}[t]
    \centering
    \IfFileExists{figura1_relaxed_pes_auditada.png}{%
        \includegraphics[width=0.88\textwidth]{figura1_relaxed_pes_auditada.png}%
    }{%
        \fbox{\parbox{0.88\textwidth}{\centering\vspace{0.9cm}\textbf{[Colocar \texttt{figura1\_relaxed\_pes\_auditada.png} en la misma carpeta]}\vspace{0.9cm}}}%
    }
    \caption{\textbf{Exploración de Nivel 1 (MACE-MP-0 + D3, \texttt{float64}, ciego al espín).} \textbf{Panel Superior:} Escaneo relajado de 13 pasos mostrando las crestas aparentes entre $2.60\text{ y }2.44\text{ \AA}$. \textbf{Panel Inferior:} Dímeros relajados en MACE-MP-0 + D3.}
    \label{fig:pes_mace}
\end{figure*}

% ==============================================================================
% FIGURA 2 A ANCHO COMPLETO (AUDITORÍA AB INITIO ORCA 6.1.1)
% ==============================================================================
\begin{figure*}[t]
    \centering
    \IfFileExists{figura2_perfiles_homogeneizados_A_y_B.png}{%
        \includegraphics[width=0.95\textwidth]{figura2_perfiles_homogeneizados_A_y_B.png}%
    }{%
        \fbox{\parbox{0.92\textwidth}{\centering\vspace{1.1cm}\textbf{[Colocar \texttt{figura2\_perfiles\_homogeneizados\_A\_y\_B.png} en la misma carpeta]}\vspace{1.1cm}}}%
    }
    \caption{\textbf{Auditoría Ab Initio Homogeneizada de Nivel 2 en ORCA 6.1.1 (\texttt{BS-UKS} $\omega\text{B97X-D3/def2-SVP}$: Opción A vs.\ Opción B vs.\ MACE-MP-0+D3).} \textbf{Paneles Superiores:} Bifurcación entre la rama repulsiva Triplete (\texttt{HS-UKS}, $\uparrow\uparrow$, roja/morada, seguida por MACE-MP-0 entre $3.40\text{ y }2.44\text{ \AA}$) y el canal reactivo Singlete (\texttt{BS-UKS} verde/naranja y proyección de Yamaguchi \texttt{AP-UKS} cian), el cual desciende sin barrera hacia el pozo covalente. \textbf{Paneles Inferiores:} Colapso continuo del valor esperado de espín $\langle \hat{S}^2 \rangle_{\text{BS}}$ desde el régimen birradicalario puro ($\langle \hat{S}^2 \rangle \approx 1.0$) en $3.40\text{ \AA}$ hasta el singlete de capa cerrada ($\langle \hat{S}^2 \rangle = 0.0$) en $\le 2.12\text{ \AA}$.}
    \label{fig:orca_audit}
\end{figure*}

% ==============================================================================
% 3. RESULTADOS: BIFURCACIÓN DE ESPÍN Y PRECISIÓN SUB-MILIANGSTROM
% ==============================================================================
\section{Resultados: Bifurcación de Espín y Pozo Covalente}

\begin{table}[H]
\centering
\caption{Comparación estructural y energética en el pozo de equilibrio ($d_{\text{eq}}$): Nivel 1 (MACE-MP-0+D3) frente a Nivel 2 (\texttt{UKS B97-3c} y $\omega\text{B97X-D3/def2-SVP}$) y difracción de rayos X ($1.578\text{ \AA}$) \cite{Alden1968}.}
\label{tab:resumen_pes}
\resizebox{\columnwidth}{!}{%
\begin{tabular}{@{}lcccccc@{}}
\toprule
\textbf{Sistema} & \textbf{$d_{\text{eq}}^{\text{MACE}}$ (\AA)} & \textbf{$d_{\text{eq}}^{\text{B97-3c}}$ (\AA)} & \textbf{Desv.\ RX (B97-3c)} & \textbf{$\Delta E_{\text{MACE}}$ (eV)} & \textbf{$\Delta E_{\text{Opt A}}$ (eV)} & \textbf{$\Delta E_{\text{Opt B}}$ (eV)} \\ \midrule
$\text{C}_{20}\text{H}_{30}$   & 1.596 & \textbf{1.5771} & $\mathbf{-0.06\%}$ & $-2.577$ & $-4.524$ & $-3.921$ \\
$\text{C}_{19}\text{H}_{29}\text{B}$ & 1.635 & \textbf{1.5810} & $+0.19\%$ & $-1.439$ & $-3.863$ & $-3.651$ \\
$\text{C}_{19}\text{H}_{29}\text{P}$ & 1.643 & \textbf{1.5880} & $+0.63\%$ & $-1.859$ & $-4.210$ & $-3.727$ \\ \bottomrule
\end{tabular}%
}
\end{table}

\subsection{Resolución del Artefacto de Barrera de MACE-MP-0}
La comparación directa entre la Figura~\ref{fig:pes_mace} y la Figura~\ref{fig:orca_audit} (resumida numéricamente en la Tabla~\ref{tab:auditoria_orca}) resuelve de forma definitiva el comportamiento del sistema entre $3.40\text{ \AA}$ y $2.44\text{ \AA}$:
\begin{enumerate}
    \item \textbf{Por qué MACE-MP-0 mostraba una cresta aparente:} Al carecer de representación explícita de espín, la superficie interpolada por MACE-MP-0 + D3 (curva azul punteada en la Figura~\ref{fig:orca_audit}) sigue entre $3.40\text{ \AA}$ y $2.44\text{ \AA}$ casi exactamente la rama del \textbf{estado Triplete ferromagnético} (\texttt{HS-UKS}, $S=1, \uparrow\uparrow$). En dicho canal triplete, el principio de exclusión de Pauli prohíbe el solapamiento enlazante de los orbitales frontales y deja expuesta la compresión estérica de los seis hidrógenos ecuatoriales de hombro, elevando la energía hasta $\Delta E_{\text{HS}} = +0.182\text{ a }+0.579\text{ eV}$ en la Opción B.
    \item \textbf{Acoplamiento Singlete Antiferromagnético Sin Barrera ($\uparrow\downarrow$):} Cuando los dos radicales se aproximan con espines antiparalelos ($M_S=0$, curvas verde y naranja en la Figura~\ref{fig:orca_audit}), el inicio del enlace covalente $\sigma(\text{C--C})$ entre $2.92\text{ \AA}$ y $2.44\text{ \AA}$ aporta una energía de estabilización de intercambio ($\Delta E_{\text{HS}} - \Delta E_{\text{BS}}$) de $0.9\text{ a }1.8\text{ eV}$ que supera por completo la repulsión estérica de los hombros. Como resultado, tanto $\Delta E_{\text{BS}}$ como el singlete puro proyectado por Yamaguchi ($\Delta E_{\text{AP}}$) **descienden de forma monótona y sin barrera de activación** desde el pre-complejo de Van der Waals ($3.40\text{ \AA}$) hasta el fondo del pozo covalente en los tres sistemas.
    \item \textbf{Transición Continua de Espín ($\langle \hat{S}^2 \rangle_{\text{BS}}$):} Los paneles inferiores de la Figura~\ref{fig:orca_audit} muestran cómo el carácter birradicalario abierto en $3.40\text{ \AA}$ ($\langle \hat{S}^2 \rangle_{\text{BS}} = 0.936\text{--}0.994$, acoplamiento magnético débil $J = -531\text{ cm}^{-1}$ en $\text{C}_{20}\text{H}_{30}$) se aparea progresivamente en $2.60\text{ \AA}$ ($\langle \hat{S}^2 \rangle_{\text{BS}} \approx 0.44\text{--}0.62$) y $2.44\text{ \AA}$ ($\langle \hat{S}^2 \rangle_{\text{BS}} \approx 0.06\text{--}0.24$), colapsando a un singlete covalente puro de capa cerrada ($\langle \hat{S}^2 \rangle_{\text{BS}} = 0.000$) a partir de $d \le 2.12\text{ \AA}$.
\end{enumerate}

\begin{table*}[t]
\centering
\caption{Auditoría cuántica completa homogeneizada en ORCA 6.1.1 (\texttt{BS-UKS} $\omega\text{B97X-D3/def2-SVP}$, 27 cálculos) frente a MACE-MP-0 + D3. Energías $\Delta E(d) = E(d) - E_\infty$ en $\text{eV}$. $\Delta E_{\text{BS}}$: singlete de simetría rota ($M_S=0, \uparrow\downarrow$); $\Delta E_{\text{AP}}$: singlete proyectado de Yamaguchi; $\Delta E_{\text{HS}}$: triplete ferromagnético ($S=1, \uparrow\uparrow$).}
\label{tab:auditoria_orca}
\resizebox{0.96\textwidth}{!}{%
\begin{tabular}{@{}lcrrrrrrrrr@{}}
\toprule
\textbf{Sistema} & \textbf{$d(\text{C--C})$ (\AA)} & \textbf{$\Delta E_{\text{MACE}}$ (eV)} & \textbf{Opt A: $\Delta E_{\text{BS}}$} & \textbf{Opt A: $\Delta E_{\text{AP}}$} & \textbf{Opt A: $\Delta E_{\text{HS}}$ ($\uparrow\uparrow$)} & \textbf{Opt A: $\langle \hat{S}^2 \rangle_{\text{BS}}$} & \textbf{Opt B: $\Delta E_{\text{BS}}$} & \textbf{Opt B: $\Delta E_{\text{AP}}$} & \textbf{Opt B: $\Delta E_{\text{HS}}$ ($\uparrow\uparrow$)} & \textbf{Opt B: $\langle \hat{S}^2 \rangle_{\text{BS}}$} \\ \midrule
$\text{C}_{20}\text{H}_{30}$ & 3.40 & $-0.333$ & $-0.247$ & $-0.366$ & $-0.116$ & 0.953 & $-0.227$ & $-0.376$ & $-0.057$ & 0.936 \\
(Intrínseco) & 2.92 & $-0.213$ & $-0.428$ & $-0.709$ & $-0.018$ & 0.816 & $-0.492$ & $-0.817$ & $+0.036$ & 0.766 \\
 & 2.60 & $-0.082$ & $-0.813$ & $-1.153$ & $+0.159$ & 0.521 & $-0.876$ & $-1.202$ & $+0.304$ & 0.436 \\
 & 2.44 & $-0.115$ & $-1.413$ & $-1.548$ & $+0.205$ & 0.155 & $-1.240$ & $-1.300$ & $+0.579$ & 0.064 \\
 & 2.12 & $-0.966$ & $-2.955$ & $-2.955$ & $+1.092$ & 0.000 & $-2.255$ & $-2.255$ & $+1.713$ & 0.000 \\
 & $d_{\text{eq}}$ ($1.596 / 1.577$) & $-2.577$ & $-4.524$ & $-4.524$ & --- & 0.000 & $\mathbf{-3.921}$ & $\mathbf{-3.921}$ & --- & 0.000 \\ \midrule
$\text{C}_{19}\text{H}_{29}\text{B}$ & 3.40 & $-0.223$ & $-0.179$ & $-0.211$ & $-0.146$ & 0.994 & $-0.215$ & $-0.323$ & $-0.097$ & 0.961 \\
(Con Boro) & 2.92 & $+0.010$ & $-0.423$ & $-0.611$ & $-0.184$ & 0.885 & $-0.425$ & $-0.691$ & $-0.042$ & 0.825 \\
 & 2.60 & $+0.333$ & $-0.845$ & $-1.165$ & $-0.127$ & 0.621 & $-0.702$ & $-1.032$ & $+0.182$ & 0.548 \\
 & 2.44 & $+0.419$ & $-1.412$ & $-1.589$ & $-0.023$ & 0.229 & $-0.954$ & $-1.139$ & $+0.423$ & 0.238 \\
 & 2.12 & $-0.156$ & $-2.673$ & $-2.673$ & $+1.075$ & 0.000 & $-1.732$ & $-1.732$ & $+1.372$ & 0.000 \\
 & $d_{\text{eq}}$ ($1.635 / 1.581$) & $-1.439$ & $-3.863$ & $-3.863$ & --- & 0.000 & $\mathbf{-3.651}$ & $\mathbf{-3.651}$ & --- & 0.000 \\ \midrule
$\text{C}_{19}\text{H}_{29}\text{P}$ & 3.40 & $-0.273$ & $-0.235$ & $-0.333$ & $-0.130$ & 0.979 & $-0.229$ & $-0.367$ & $-0.074$ & 0.942 \\
(Con Fósforo) & 2.92 & $-0.085$ & $-0.521$ & $-0.779$ & $-0.159$ & 0.843 & $-0.477$ & $-0.785$ & $+0.011$ & 0.779 \\
 & 2.60 & $+0.127$ & $-1.066$ & $-1.393$ & $-0.120$ & 0.519 & $-0.825$ & $-1.148$ & $+0.277$ & 0.456 \\
 & 2.44 & $+0.139$ & $-1.706$ & $-1.809$ & $-0.075$ & 0.119 & $-1.152$ & $-1.235$ & $+0.558$ & 0.093 \\
 & 2.12 & $-0.549$ & $-3.057$ & $-3.057$ & $+1.003$ & 0.000 & $-2.058$ & $-2.058$ & $+1.726$ & 0.000 \\
 & $d_{\text{eq}}$ ($1.643 / 1.588$) & $-1.859$ & $-4.210$ & $-4.210$ & --- & 0.000 & $\mathbf{-3.727}$ & $\mathbf{-3.727}$ & --- & 0.000 \\ \bottomrule
\end{tabular}%
}
\end{table*}

\subsection{Exactitud Cristalográfica y Corrección Trans-Jaula en Boro}
La optimización completa con \texttt{UKS B97-3c} aporta dos validaciones estructurales críticas (Tabla~\ref{tab:resumen_pes}):
\begin{itemize}
    \item \textbf{Contraste con Rayos X a Nivel Sub-miliangstrom:} En el dímero intrínseco ($\text{C}_{20}\text{H}_{30}$), \texttt{UKS B97-3c} predice una longitud de enlace central $d_{\text{eq}}(\text{C--C}) = 1.5771\text{ \AA}$, reproduciendo el valor experimental por difracción de rayos X de monocristal del 1,1'-biadamantilo ($1.578\text{ \AA}$, Alden et al.\ \cite{Alden1968}) con una desviación de apenas $\mathbf{-0.06\%}$ ($-0.0009\text{ \AA}$), mejorando sustancialmente el $+1.14\%$ ($1.596\text{ \AA}$) de MACE-MP-0.
    \item \textbf{Corrección de la Sobre-Contracción del Boro:} En el radical monomérico sustituido con boro ($\text{C}_9\text{BH}_{14}^\bullet$), MACE-MP-0 sobreestimaba la contracción trans-jaula entre el boro trivalente y el carbono radical ($1.777\text{ \AA}$), lo que inflaba artificialmente la penalización de reorganización del dímero ($\Delta E_{\text{MACE}} = -1.439\text{ eV}$). Al optimizar con \texttt{UKS B97-3c}, la distancia trans-jaula en el monómero se estabiliza en $2.149\text{ \AA}$ y el hetero-dímero $\text{C}_{19}\text{H}_{29}\text{B}$ alcanza un pozo covalente profundo de $\Delta E_{\text{Opt B}} = -3.651\text{ eV}$ ($d_{\text{eq}} = 1.581\text{ \AA}$), muy cercano al derivado con fósforo ($\Delta E_{\text{Opt B}} = -3.727\text{ eV}$, $d_{\text{eq}} = 1.588\text{ \AA}$) y al sistema intrínseco ($-3.921\text{ eV}$).
\end{itemize}

\subsection{Estadística Cuántica de Espín ($1:3$) en un Reactor Real}
La bifurcación entre las ramas \texttt{BS-UKS} ($\uparrow\downarrow$) y \texttt{HS-UKS} ($\uparrow\uparrow$) impone una restricción física fundamental sobre la cinética de colisión en vacío. En un haz molecular térmico no polarizado magnéticamente, el acoplamiento de dos espines $S=1/2$ genera cuatro estados degenerados en ausencia de campo: un estado \textbf{Singlete} reactivo ($S=0$: $\frac{1}{\sqrt{2}}(\lvert\uparrow\downarrow\rangle - \lvert\downarrow\uparrow\rangle)$) y tres estados \textbf{Triplete} ($S=1$: $\lvert\uparrow\uparrow\rangle, \lvert\downarrow\downarrow\rangle, \frac{1}{\sqrt{2}}(\lvert\uparrow\downarrow\rangle + \lvert\downarrow\uparrow\rangle)$).

Por consiguiente, no debe afirmarse que toda colisión entre radicales sea libre de barrera: únicamente el $\mathbf{25\%}$ ($1/4$) de los encuentros bimoleculares ingresa por el canal singlete antiferromagnético ($\uparrow\downarrow$) y cae sin barrera hacia el pozo covalente ($\Delta E = -3.65\text{ a }-3.92\text{ eV}$), mientras que el $\mathbf{75\%}$ ($3/4$) restante ingresa por el canal triplete ($\uparrow\uparrow$) y choca contra una pared repulsiva de Pauli + estérica ($\Delta E_{\text{HS}} > 0\text{ eV}$ para $d \le 2.9\text{ \AA}$), sufriendo rebote térmico elástico a menos que ocurra cruce entre sistemas (\textit{intersystem crossing}, ISC) asistido por superficie o pre-polarización de espín.

% ==============================================================================
% 4. MATRIZ DE HIPÓTESIS, REFUTABILIDAD Y REACTOR CONCEPTUAL
% ==============================================================================
\section{Alcance, Hipótesis y Criterios de Refutación}
Aun con la validación \textit{ab initio} en ORCA 6.1.1 completada, \textbf{estos resultados acotan exclusivamente la dimerización bimolecular aislada y todavía no demuestran la síntesis de una película continua de diamante, ni la propagación multicapa, ni la activación eléctrica del dopaje semiconductor}:
\begin{itemize}\setlength{\itemsep}{1pt}\setlength{\parskip}{0pt}
    \item \textbf{$H_1$ (Dimerización Singlete Sin Barrera -- Validada):} A nivel \texttt{BS-UKS} $\omega\text{B97X-D3/def2-SVP}$ // \texttt{UKS B97-3c}, el acoplamiento coaxial de dos radicales diamantoides carece de barrera en el canal singlete ($25\%$ estadístico), con enlaces inter-jaula de $1.577\text{--}1.588\text{ \AA}$ y energías de $-3.65\text{ a }-3.92\text{ eV}$.
    \item \textbf{$H_2$ y $H_3$ (Hipótesis Pendientes):} (i) \textit{Activación $H_2$:} requiere demostrar abstracción selectiva $\text{C--H}$ ($\sim 4.3\text{ eV}$) mediante electrones suaves ($0.5\text{--}2.0\text{ keV}$) o VUV ($172\text{ nm}$) sin fragmentar la red $\text{C--C}$ \cite{Fu2025}. (ii) \textit{Propagación $H_3$:} requiere modelar el crecimiento 3D en sustrato donde el $75\%$ triplete relaje su espín sin desorberse y evaluar la conductividad activa.
    \item \textbf{Criterios de Refutación:} El modelo se refutaría si la fragmentación $\text{C--C}$ domina sobre la abstracción $\text{C--H}$, o si la tasa de desorción térmica en vacío impide superar el límite estadístico de espín ($1:3$).
\end{itemize}

\begin{figure}[H]
    \centering
    \IfFileExists{figura2_reactor_tharos.png}{%
        \includegraphics[width=\columnwidth]{figura2_reactor_tharos.png}%
    }{%
        \fbox{\parbox{0.95\columnwidth}{\centering\small\textbf{[Colocar \texttt{figura2\_reactor\_tharos.png} en la carpeta]}}}%
    }
    \caption{\textbf{Esquema conceptual del banco de alto vacío propuesto para someter a prueba $H_2$ y $H_3$.} Co-sublimación desde celdas Knudsen ($60\text{--}80\text{ }^\circ\text{C}$), excitación electrónica ($0.5\text{--}2.0\text{ keV}$) o VUV ($172\text{ nm}$) y sustrato termorregulado ($150\text{--}250\text{ }^\circ\text{C}$).}
    \label{fig:reactor_tharos}
\end{figure}

Únicamente como marco experimental futuro para someter a prueba $H_2$ y $H_3$, la Figura~\ref{fig:reactor_tharos} ilustra la configuración conceptual donde vapores dosificados desde celdas Knudsen en alto vacío ($\sim 10^{-6}\text{ Torr} / 10^{-4}\text{ Pa}$) incidirían sobre un sustrato a $150\text{--}250\text{ }^\circ\text{C}$ bajo excitación electrónica suave o VUV.

% ==============================================================================
% 5. CONCLUSIÓN Y DECLARACIONES ACADÉMICAS
% ==============================================================================
\section{Conclusión}
La auditoría cuántica de espín abierto en ORCA 6.1.1 (\texttt{BS-UKS} $\omega\text{B97X-D3/def2-SVP}$ // \texttt{UKS B97-3c}) demuestra que la aparente barrera estérica observada inicialmente con el potencial neuronal MACE-MP-0 + D3 ($0.251\text{--}0.642\text{ eV}$) correspondía a la rama repulsiva del estado Triplete ($S=1, \uparrow\uparrow$). En el canal Singlete Antiferromagnético ($M_S=0, \uparrow\downarrow$, que representa el $25\%$ estadístico de las colisiones térmicas no polarizadas), la estabilización de intercambio supera la compresión de los hidrógenos ecuatoriales, permitiendo una dimerización sin barrera hacia pozos covalentes de $-3.65\text{ a }-3.92\text{ eV}$ y reproduciendo a $-0.06\%$ ($1.5771\text{ \AA}$) la estructura cristalográfica por rayos X del 1,1'-biadamantilo. Estos hallazgos establecen una base cuántica rigurosa para el paso elemental de dimerización en vacío, dejando abiertos los desafíos de activación selectiva $\text{C--H}$ y propagación cristalina en superficie.

\section*{Data and Code Availability}
\footnotesize
El paquete de auditoría computacional (\texttt{paquete\_auditoria\_priones.zip}) y la carpeta de validación \textit{ab initio} (\texttt{dft\_orca\_validation}), que incluyen los scripts de simulación, los 27 archivos de entrada/salida de ORCA 6.1.1 (\texttt{.inp}, \texttt{.out}, \texttt{.xyz}), la tabla consolidada (\texttt{tabla\_auditoria\_completa\_mace\_vs\_orca\_A\_y\_B.csv}) y el registro JSON (\texttt{auditoria\_completa\_5puntos.json}), se encuentran depositados bajo licencia MIT en Zenodo (\textbf{DOI:} \href{https://doi.org/10.5281/zenodo.23427695}{\texttt{10.5281/zenodo.23427695}}) y vinculados al perfil de código del autor (\texttt{NocWoeder165}).

\section*{CRediT Authorship Contribution Statement}
\footnotesize
\textbf{César Leonardo Rodríguez García} \textit{(Dark Horse / NocWoeder165)}: Conceptualization, Methodology, Software, Validation, Formal analysis, Investigation, Data curation, Writing -- original draft, Writing -- review \& editing, Visualization.

\section*{Declaration of Competing Interest}
\footnotesize
Se declara la ausencia de intereses financieros en competencia o relaciones personales que pudieran haber influido en los resultados reportados en este pre-print.

% ==============================================================================
% REFERENCIAS BIBLIOGRÁFICAS (LAS 15 REFERENCIAS COMPLETAS Y VERIFICADAS)
% ==============================================================================
\begin{thebibliography}{15}
\footnotesize

\bibitem{Koizumi2001}
S.~Koizumi, K.~Watanabe, M.~Hasegawa, H.~Kanda, ``Ultraviolet emission from a diamond pn junction,'' \textit{Science} \textbf{292} (5523) (2001) 1899--1901. doi:10.1126/science.1060258.

\bibitem{Suzuki2024}
Y.~Suzuki, et al., ``Raman study on the pressure-induced phase transformation of nanographite at room temperature,'' \textit{Carbon} \textbf{225} (2024) 119075. doi:10.1016/j.carbon.2024.119075.

\bibitem{Kamo1983}
M.~Kamo, Y.~Sato, S.~Matsumoto, N.~Setaka, ``Diamond synthesis from gas phase in microwave plasma,'' \textit{J.\ Cryst.\ Growth}\ \textbf{62} (3) (1983) 642--644. doi:10.1016/0022-0248(83)90411-6.

\bibitem{Schleyer1957}
P.~v.~R.~Schleyer, ``A simple preparation of adamantane,'' \textit{J.\ Am.\ Chem.\ Soc.}\ \textbf{79} (1957) 3292. doi:10.1021/ja01569a086.

\bibitem{Dahl2003}
J.E.~Dahl, S.G.~Liu, R.M.K.~Carlson, ``Isolation and structural characterization of the higher diamondoids,'' \textit{Science} \textbf{299} (2003) 96--99. doi:10.1126/science.1078239.

\bibitem{Yang2007}
W.L.~Yang, J.D.~Fabbri, T.M.~Willey, J.R.I.~Lee, J.E.~Dahl, R.M.K.~Carlson, P.R.~Schreiner, A.A.~Fokin, B.A.~Tkachenko, N.A.~Fokina, W.~Meevasana, N.~Mannella, K.~Tanaka, X.J.~Zhou, T.~van~Buuren, M.A.~Kelly, Z.~Hussain, N.A.~Melosh, Z.-X.~Shen, ``Monochromatic electron photoemission from diamondoid monolayers,'' \textit{Science} \textbf{316} (5830) (2007) 1460--1462. doi:10.1126/science.1141811.

\bibitem{Fu2025}
J.~Fu, T.~Nakamuro, E.~Nakamura, ``Nanodiamond synthesis from adamantane by electron-beam driven C--H activation at 100 K,'' \textit{Science} \textbf{389} (2025) eadw2025 (PMID: 40906847). doi:10.1126/science.adw2025.

\bibitem{Gebbie2018}
M.A.~Gebbie, H.~Ishiwata, P.J.~McQuade, J.E.~Dahl, P.R.~Schreiner, Z.-X.~Shen, N.A.~Melosh, et al., ``Experimental measurement of the diamond nucleation landscape reveals classical and nonclassical features,'' \textit{Proc.\ Natl.\ Acad.\ Sci.\ U.S.A.}\ \textbf{115} (33) (2018) 8284--8289. doi:10.1073/pnas.1803654115.

\bibitem{Batatia2024}
I.~Batatia, et al., ``A foundation model for atomistic materials chemistry,'' \textit{arXiv:2401.00096} (2024). doi:10.48550/arXiv.2401.00096.

\bibitem{Grimme2010}
S.~Grimme, J.~Antony, S.~Ehrlich, H.~Krieg, ``A consistent and accurate ab initio parametrization of density functional dispersion correction (DFT-D3),'' \textit{J.\ Chem.\ Phys.}\ \textbf{132} (2010) 154104. doi:10.1063/1.3382344.

\bibitem{Neese2022}
F.~Neese, ``Software update: The ORCA program system---Version 5.0,'' \textit{WIREs Comput.\ Mol.\ Sci.}\ \textbf{12} (5) (2022) e1606. doi:10.1002/wcms.1606.

\bibitem{Chai2013}
Y.-S.~Lin, G.-D.~Li, S.-P.~Mao, J.-D.~Chai, ``Long-range corrected hybrid density functionals with improved dispersion corrections,'' \textit{J.\ Chem.\ Theory Comput.}\ \textbf{9} (1) (2013) 263--272. doi:10.1021/ct300715s.

\bibitem{Brandenburg2018}
J.G.~Brandenburg, C.~Bannwarth, A.~Hansen, S.~Grimme, ``B97-3c: A revised low-cost variant of the B97-D density functional method,'' \textit{J.\ Chem.\ Phys.}\ \textbf{148} (6) (2018) 064104. doi:10.1063/1.5012601.

\bibitem{Yamaguchi1988}
K.~Yamaguchi, F.~Jensen, A.~Dorigo, K.N.~Houk, ``A spin correction procedure for unrestricted Hartree-Fock and M\o ller-Plesset wavefunctions for singlet diradicals and polyradicals,'' \textit{Chem.\ Phys.\ Lett.}\ \textbf{149} (5--6) (1988) 537--542. doi:10.1016/0009-2614(88)80378-6.

\bibitem{Alden1968}
R.A.~Alden, J.~Kraut, T.G.~Traylor, ``Molecular and crystal structure of 1,1'-biadamantyl,'' \textit{J.\ Am.\ Chem.\ Soc.}\ \textbf{90} (1) (1968) 74--82. doi:10.1021/ja01003a014.

\end{thebibliography}

\end{document}
